\subsection{共轭性的应用}

定理 2. 设 \(\mathfrak{g}\) 是一个李代数。

\begin{enumerate}
\def\labelenumi{(\roman{enumi})}
\item
  \(\mathfrak{g}\) 的诸嘉当子代数都有相同的维数，即 \(\operatorname{rk}(\mathfrak{g})\) ，以及相同的幂零类。
\item
  元素 \(x \in \mathfrak{g}\) 正则当且仅当 \(\mathfrak{g}^0(x)\) 是 \(\mathfrak{g}\) 的嘉当子代数；每个嘉当子代数都这样得到。
\end{enumerate}

为证 (i)，可以假定 k 代数闭（参看~§2，命题 3 和命题 6），这时它由第 2 节的定理 1 得出。断言 (ii) 由 (i) 以及 §2 定理 1 的 (i) 和 (iv) 得出。

命题 3. 设 \(\mathfrak{g}\) 是一个李代数，\(\mathfrak{g}'\) 是 \(\mathfrak{g}\) 的一个子代数。下列条件等价：

\begin{enumerate}
\def\labelenumi{(\roman{enumi})}
\item
  \(\mathfrak{g}'\) 包含 \(\mathfrak{g}\) 的一个正则元，且 \(\mathrm{rk}(\mathfrak{g}) = \mathrm{rk}(\mathfrak{g}')\) ；
\item
  \(\mathfrak{g}'\) 包含 \(\mathfrak{g}\) 的一个嘉当子代数；
\item
  \(\mathfrak{g}'\) 的每个嘉当子代数都是 \(\mathfrak{g}\) 的嘉当子代数。
\item
  \(\Longrightarrow\) (ii)：假定 \(\operatorname{rk}(\mathfrak{g}) = \operatorname{rk}(\mathfrak{g}')\) ，且存在 \(x \in \mathfrak{g}'\) 在 \(\mathfrak{g}\) 中正则。设 \(\mathfrak{h} = \mathfrak{g}^0(x)\) ，\(\mathfrak{h}' = \mathfrak{g}'^0(x) = \mathfrak{h} \cap \mathfrak{g}'\) 。那么
\end{enumerate}

\[
\operatorname{rk} \left(\mathfrak {g} ^ {\prime}\right) \leq \dim \mathfrak {h} ^ {\prime} \leq \dim \mathfrak {h} = \operatorname{rk} (\mathfrak {g}) = \operatorname{rk} \left(\mathfrak {g} ^ {\prime}\right)
\]

故 \(\mathfrak{h} = \mathfrak{h}'\subset \mathfrak{g}'\) 。这就证明了 (ii)。

\begin{enumerate}
\def\labelenumi{(\roman{enumi})}
\setcounter{enumi}{1}
\item
  \(\Longrightarrow\) (iii)：假定 \(g'\) 包含 g 的一个嘉当子代数 h，并设 \(h_{1}\) 是 \(g'\) 的一个嘉当子代数。为证 \(h_{1}\) 是 g 的嘉当子代数，可以假定 k 代数闭。设 \(\mathrm{E}(\mathfrak{h})\) 和 \(\mathrm{E}'(\mathfrak{h})\) 是与 h 相关联的 g 和 \(g'\) 的自同构群（第 2 节）。由定理 1，存在 \(f \in \mathrm{E}'(\mathfrak{h})\) 使 \(f(\mathfrak{h}) = \mathfrak{h}_{1}\) 。而 \(\mathrm{E}'(\mathfrak{h})\) 的每个元素都由 \(\mathrm{E}(\mathfrak{h})\) 的一个元素诱导；事实上，只需对 \(e^{\operatorname{ad} x}\) （其中 \(x \in \mathfrak{g}'^{\lambda}(\mathfrak{h}), \lambda \neq 0\) ）验证这一点，这时它由包含关系 \(\mathfrak{g}'^{\lambda}(\mathfrak{h}) \subset \mathfrak{g}^{\lambda}(\mathfrak{h})\) 得出。于是 \(h_{1}\) 是 g 的嘉当子代数。
\item
  \(\Longrightarrow\) (i)：假定条件 (iii) 满足。设 h 是 \(g'\) 的一个嘉当子代数。由于它是 g 的嘉当子代数，它包含 g 的一个正则元（定理 2 (ii)），另一方面 \(\operatorname{rk}(\mathfrak{g}) = \dim(\mathfrak{h}) = \operatorname{rk}(\mathfrak{g}')\) 。
\end{enumerate}

推论. 设 \(\mathfrak{h}\) 是 \(\mathfrak{g}\) 的一个幂零子代数。子代数 \(\mathfrak{g}^0 (\mathfrak{h})\) 具有命题 3 中的性质 (i)、(ii)、(iii)。

事实上，§2 第 3 节的命题 11 表明 \(\mathfrak{g}^0 (\mathfrak{h})\) 具有性质 (ii)。

命题 4. 设 g 是一个李代数，l 是 g 的秩，c 是 g 的嘉当子代数的幂零类，\(x \in g\) 。存在 g 的一个维数为 l、幂零类 \(\leq c\) 且包含 x 的子代数。

设 T 是一个不定元。设 \(k' = k(T)\) ，\(g' = g \otimes_k k'\) 。若 h 是 g 的嘉当子代数，则 \(h \otimes_k k'\) 是 \(g'\) 的嘉当子代数，故 \(g'\) 的秩是 l，且 \(g'\) 的嘉当子代数的幂零类是 c。

选取正则元 \(y\) （\(\mathfrak{g}\) 的）。用 §2 第 2 节的记号，我们有 \(a_l(y) \neq 0\) 。仍用 \(a_l\) 记 \(\mathfrak{g}'\) 上延拓 \(a_l\) 的多项式函数。那么元素 \(a_l(x + \mathrm{T}y)\) （\(k[\mathrm{T}]\) 的）的支配系数是 \(a_l(y)\) 。特别地，\(x + \mathrm{T}y\) 在 \(\mathfrak{g}'\) 中正则。设 \(\mathfrak{h}'\) 是 \(\operatorname{ad}(x + \mathrm{T}y)\) 在 \(\mathfrak{g}'\) 中的幂零空间。那么 \(\dim \mathfrak{h}' = l\) ，且 \(\mathfrak{h}'\) 的幂零类是 \(c\) 。设 \(\mathfrak{k} = \mathfrak{h}' \cap (\mathfrak{g} \otimes_k k[\mathrm{T}])\) ；那么 \(\mathfrak{k} \otimes_{k[\mathrm{T}]} k(\mathrm{T}) = \mathfrak{h}'\) 。

设 \(\varphi\) 是从 \(k[\mathrm{T}]\) 到 \(k\) 的使 \(\varphi (\mathrm{T}) = 0\) 的同态，并设 \(\psi\) 是同态 \(1\otimes \varphi\) （从 \(\mathfrak{g}\otimes_k k[\mathrm{T}]\) 到 \(\mathfrak{g}\) ）。那么 \(\psi (\mathfrak{k})\) 是 \(\mathfrak{g}\) 的一个幂零类 \(\leq c\) 且包含 \(\psi (x + \mathrm{T}y) = x\) 的子代数。

在自由 \(k[\mathrm{T}]\) -模 \(\mathfrak{g} \otimes_k k[\mathrm{T}]\) 中，\(\mathfrak{k}\) 是秩为 \(l\) 的子模，且 \((\mathfrak{g} \otimes_k k[\mathrm{T}]) / \mathfrak{k}\) 无挠，故子模 \(\mathfrak{k}\) 是 \(\mathfrak{g} \otimes_k k[\mathrm{T}]\) 的直和项（《代数》，第七章，§4，第 2 节，定理 1）。因此 \(\dim_k \psi(\mathfrak{k}) = l\) ，证毕。
% semantic-fragments: legacy-input-seam
\relax{}
